\section{Forelesning 2: \emph{Strategisk analyse (del I) og typiske strategiske virkemidler}}
\label{sec:lec02}

\subsection*{Oversikt}
Forelesning 2 introduserer verktøykassen for \emph{ekstern} og \emph{intern}
strategisk analyse, og en katalog av typiske strategiske grep en aktør kan velge.
Målet er å gå fra Porters abstrakte posisjoneringsbegrep (\autoref{sec:lec01})
til konkrete analyseverktøy som strukturerer hva man ser etter når man
diagnostiserer en virksomhets situasjon.

\subsection{Hovedteori}

\subsubsection*{Ekstern analyse: Porters Five Forces}
Rammeverket kartlegger konkurransedynamikken i en bransje gjennom fem krefter:
\begin{enumerate}
  \item \term{Industry rivalry} -- intensiteten av konkurranse mellom
    eksisterende aktører.
  \item \term{Threat of new entrants} -- inntrengertrussel, knyttet til
    inngangsbarrierer.
  \item \term{Bargaining power of buyers} -- kjøpermakt.
  \item \term{Bargaining power of suppliers} -- leverandørmakt.
  \item \term{Threat of substitutes} -- substituttrusler.
\end{enumerate}
Den samlede styrken av de fem kreftene bestemmer bransjens strukturelle
lønnsomhetspotensial.

\begin{definition}{Bransjeattraktivitet}
En bransje er strukturelt attraktiv når de fem kreftene er svake. Strategisk
posisjon handler om å plassere seg slik at de fem kreftene treffer en mindre
hardt enn konkurrentene.
\end{definition}

\subsubsection*{Intern analyse: Verdikjeden}
Porters \term{value chain} bryter virksomhetens aktiviteter ned i:
\begin{itemize}
  \item \emph{Primæraktiviteter}: inngående logistikk, drift, utgående
    logistikk, markedsføring/salg, service.
  \item \emph{Støtteaktiviteter}: virksomhetens infrastruktur, HRM,
    teknologiutvikling, innkjøp.
\end{itemize}
Hver aktivitet er en mulig kilde til konkurransefortrinn -- enten gjennom lavere
kostnad eller gjennom differensiering.

\subsubsection*{Ressursbasert syn (kort)}
Komplementært til Porters posisjoneringsperspektiv: noen virksomheter har
\term{VRIN-ressurser} (verdifulle, sjeldne, uimiterbare, ikke-substituerbare)
som gir varig fortrinn. Relevant for å vurdere f.eks.\ Coops medlemsdata.

\subsection{Sentrale fagbegreper}
\begin{itemize}
  \item \term{Five Forces} -- struktur for ekstern bransjeanalyse.
  \item \term{Value chain} -- struktur for intern aktivitetsanalyse.
  \item \term{Entry barriers} -- skalafordeler, brand, byttekostnader, kapital.
  \item \term{Switching costs} -- kunders kostnad ved å bytte leverandør.
  \item \term{Network effects} -- verdi som vokser med antall brukere.
  \item \term{Lock-in} -- mekanisme som binder kunder til en leverandør.
\end{itemize}

\subsection{Modeller og rammeverk}

\figureplaceholder{lec02-porter-five-forces}{Porters Five Forces -- de fem
konkurransekreftene som bestemmer bransjelønnsomhet.}

\figureplaceholder{lec02-value-chain}{Porters verdikjede -- primær- og
støtteaktiviteter.}

\subsubsection*{Typiske strategiske virkemidler}
For en etablert aktør som møter digital disrupsjon (jf.\ \autoref{sec:lec05}):
\begin{enumerate}
  \item \textbf{Forsvare} eksisterende posisjon (styrke eksisterende fortrinn).
  \item \textbf{Investere} i nye digitale kanaler.
  \item \textbf{Kjøpe} digitale kapabiliteter (oppkjøp).
  \item \textbf{Samarbeide} med digitale aktører (allianser, plattformer).
  \item \textbf{Selv-disrupsjon} -- bygge en disruptiv enhet før andre gjør det.
\end{enumerate}
Valget avhenger av strategisk fit med forretningsmodellen (\autoref{sec:lec03}).

\subsection{Eksempler og caser}
\begin{itemize}
  \item \textbf{Dagligvare i Danmark} -- Five Forces analyse: høy rivaliditet
    (Salling, Coop, REMA, Lidl), moderat trussel fra netthandel (Nemlig),
    sterk leverandørmakt fra store merker, lav kjøpermakt per kunde men
    sterk samlet, substitutter inkluderer ferdigmat-bud.
  \item \textbf{Coops verdikjede} -- innkjøp og logistikk er sentralt;
    digital teknologi i butikk (Scan \& Betal, lojalitetsapp) hører hjemme
    under \emph{teknologiutvikling} og \emph{markedsføring}.
\end{itemize}

\subsection{Eksamensrelevans}
\begin{itemize}
  \item Kunne gjennomføre en Five Forces-analyse på en konkret bransje.
  \item Kunne plassere digitale initiativ i verdikjeden og forklare hvor
    fortrinn oppstår.
  \item Kjenne katalogen av strategiske virkemidler og kunne velge mellom dem.
  \item Forstå at ekstern + intern analyse er to nødvendige sider av samme sak.
\end{itemize}
